%
% Arrows
%

\newcommand{\multirightarrow}[1]{%
	\mathrel{%
		\substack{%
			\let\multiarrowcontent\empty
			\foreach \i in {1,...,#1} {
	    		\ifnum\i=#1
	    			\expandafter\gappto\expandafter\multiarrowcontent\expandafter{\longrightarrow}
	    		\else
	    			\expandafter\gappto\expandafter\multiarrowcontent\expandafter{\longrightarrow\\[-0.6ex]}
	    		\fi
			} 
			\multiarrowcontent%
		}%
	}%
}

\newcommand{\multileftarrow}[1]{%
	\mathrel{%
		\substack{%
			\let\multiarrowcontent\empty
			\foreach \i in {1,...,#1} {
	    		\ifnum\i=#1
	    			\expandafter\gappto\expandafter\multiarrowcontent\expandafter{\longleftarrow}
	    		\else
	    			\expandafter\gappto\expandafter\multiarrowcontent\expandafter{\longleftarrow\\[-0.6ex]}
	    		\fi
			} 
			\multiarrowcontent%
		}%
	}%
}

\renewcommand{\twoheadrightarrow}{\mathrel{\mathrlap{\rightarrow}\mkern.5mu\rightarrow}}
\DeclareRobustCommand\longtwoheadrightarrow{\relbar\joinrel\twoheadrightarrow}
\renewcommand{\twoheadleftarrow}{\mathrel{\mathrlap{\leftarrow}\mkern.5mu\leftarrow}}
\DeclareRobustCommand\longtwoheadleftarrow{\twoheadleftarrow\joinrel\relbar}

%
% Derivatives
%

\newcommand{\dder}[2][]{%
	\ifthenelse{\equal{#1}{}}{%
		\frac{\mathrm{d}}{\mathrm{d} #2}%
	}{%
		\frac{\mathrm{d} #1}{\mathrm{d} #2}%
	}%
}
\newcommand{\parder}[2][]{%
	\ifthenelse{\equal{#1}{}}{%
		\frac{\partial}{\partial #2}%
	}{%
		\frac{\partial #1}{\partial #2}%
	}%
}
\newcommand{\delder}[2][]{%
	\ifthenelse{\equal{#1}{}}{%
		\frac{\delta}{\delta #2}%
	}{%
		\frac{\delta #1}{\delta #2}%
	}%
}

%
% Inner products
%

\newcommand{\inner}[2]{\langle#1,#2\rangle}
\newcommand{\innerLarge}[2]{\left\langle#1,#2\right\rangle}

%
% Vectors
%

\newcommand{\colvec}[1]{%
    \begingroup
        \newcommand\listcontent{}
        \def\do##1{\appto\listcontent{##1\\}}
        \expandafter\docsvlist\expandafter{#1}
        \begin{pmatrix}
            \listcontent
        \end{pmatrix}
    \endgroup
}

%
% Integrals
%

\DeclareMathOperator*{\sumint}{%
\mathchoice%
	{\ooalign{$\displaystyle\sum$\cr\hidewidth$\displaystyle\int$\hidewidth\cr}}
	{\ooalign{\raisebox{.14\height}{\scalebox{.7}{$\textstyle\sum$}}\cr\hidewidth$\textstyle\int$\hidewidth\cr}}
	{\ooalign{\raisebox{.2\height}{\scalebox{.6}{$\scriptstyle\sum$}}\cr$\scriptstyle\int$\cr}}
	{\ooalign{\raisebox{.2\height}{\scalebox{.6}{$\scriptstyle\sum$}}\cr$\scriptstyle\int$\cr}}
}

\newcommand{\cpv}{\setbox0=\hbox{$\int$}\int\hskip-\wd0{}\hskip-\wd0{~-~}} % Cauchy principal value

%
% Miscellaneous
%

\newcommand{\acton}{\vartriangleright} % action
\newcommand{\Ad}{\mathrm{Ad}} % adjoint map
\newcommand{\ad}{\mathrm{ad}} % adjoint map
\newcommand{\diag}{\mathrm{diag}} % diagonal matrix
\newcommand{\lvac}{\langle 0|} % vacuum
\newcommand{\dual}{^\vee} % vector space dual
\newcommand{\id}{\mathrm{id}} % identity
\newcommand{\sfid}{\mathsf{id}} % identity
\newcommand{\im}{\mathrm{im}} % image
\renewcommand{\Im}{\mathfrak{Im}} % imaginary part
\newcommand{\pr}{\mathsf{pr}} % projection
\renewcommand{\Re}{\mathfrak{Re}} % real part
\newcommand{\spacer}[1]{\vrule width0pt height0pt depth#1\relax}
\newcommand{\spn}{\mathrm{span}} % span
\newcommand{\str}{\mathrm{str}} % super trace
\newcommand{\T}{^{\mathrm{T}}} % transpose
\newcommand{\tocut}[1]{{\color{green} #1 }}
\newcommand{\todo}[1]{{\color{red} #1 }}
\newcommand{\tr}{\mathrm{tr}} % trace
\newcommand{\rvac}{|0\rangle} % vacuum
\newcommand{\uline}[1]{\ul{\emph{#1}}}
\newcommand{\unit}{\mathbbm{1}} % unit matrix
\newcommand{\rint}{\mathrm{int}}
\newcommand{\rtree}{\mathrm{tree}}
\newcommand{\lap}{\Delta}
\newcommand{\wave}{\mathop\Box}
\newcommand{\intprod}{\mathbin{\raisebox{\depth}{\scalebox{1}[-1]{$\lnot$}}}} % interior product

%
% Prepositions
%

\newcommand{\efor}{{~~~\mbox{for}~~~}}
\newcommand{\eforall}{{~~~\mbox{for all}~~~}}
\newcommand{\eand}{{~~~\mbox{and}~~~}}
\newcommand{\eon}{{~~\mbox{on}~~}}
\newcommand{\eor}{{~~~\mbox{or}~~~}}  		
\newcommand{\ewith}{{~~~\mbox{with}~~~}}

%
% Young diagrams
%

\newcommand{\tyng}[1]{\hbox{\tiny$\yng #1$}} % small Young diagram
\newcommand{\tyoung}[1]{\hbox{\tiny$\young #1$}} % small Young diagram